% 초안: supplementary.tex 의 Remark rem:theta-TU 바로 뒤 (Proof of Lemma phi 절) 에 넣을 후보.
% 표기는 원고 그대로: follower 는 d_F^T y 최소화, 리더 평가 phi = max d_L^T y (follower 최적해 중), Y = {y >= 0 : Ay <= r},
% z_F(r) = min d_F over Y, L(theta) = max_{y in Y} [ d_L^T y + theta (z_F(r) - d_F^T y) ] (Lemma phi (ii) 증명의 L).

\begin{remark}[Circuit bound on the multiplier]\label{rem:theta-circuit}
The constant $\theta^{U}$ enters the reformulation only through the
fixed-$\theta$ form of~\eqref{eq:phi-min} in Lemma~\ref{lem:phi}(ii), that is,
through the identity $L(\theta^{U}) = \phi(x,h,\xi)$ for every right-hand side
$r$. Any constant with this property may therefore replace $\theta^{U}$
throughout, and the following one is often much smaller. Call
$g \in \mathbb{R}^{n_y}\setminus\{0\}$ a \emph{circuit} of $A$ if the support
of $(Ag, g)$ is inclusion-minimal among nonzero vectors of this form, and let
\[
    \theta_{C} \;:=\; \max\Big\{ \frac{(d_L^\top g)^{+}}{d_F^\top g} \;:\; g \text{ a circuit of } A,\ d_F^\top g > 0 \Big\}.
\]
Circuits depend only on $A$ and are finitely many up to scaling, so $\theta_{C}$
is a uniform constant. We claim that $L(\theta) = \phi(x,h,\xi)$ for every
$\theta \ge \theta_{C}$ and every $r$.

Fix $r$. By Assumption~\ref{asm:RCR}, $\mathcal{Y} := \{y \ge 0 : Ay \le r\}$
is a polytope. Let $F$ be the face of minimizers of $d_F^\top y$ over
$\mathcal{Y}$, and let $y^{*}$ be a vertex of $\mathcal{Y}$ in $F$ that
maximizes $d_L^\top y$ over $F$, so that $\phi(x,h,\xi) = d_L^\top y^{*}$.
Every $y \in \mathcal{Y}$ satisfies $y - y^{*} = \sum_{k}\lambda_k g_k$ with
$\lambda_k \ge 0$, where the $g_k$ are the directions of the edges of
$\mathcal{Y}$ at $y^{*}$. Each edge is defined by $n_y - 1$ linearly
independent rows of $[A;\,-I]$ that are active along it, and its direction
spans the null space of these rows; any nonzero $g'$ for which the support of
$(Ag', g')$ is contained in that of $(Ag_k, g_k)$ vanishes on the same rows and
is a multiple of $g_k$, so $g_k$ is a circuit. Since $y^{*}$ minimizes
$d_F^\top y$, $d_F^\top g_k \ge 0$. If $d_F^\top g_k = 0$, the edge lies in
$F$ and $d_L^\top g_k \le 0$ by the choice of $y^{*}$; otherwise
$d_L^\top g_k \le \theta_{C}\, d_F^\top g_k$. Hence
\[
    d_L^\top (y - y^{*}) \;\le\; \theta_{C}\, d_F^\top (y - y^{*}) \;=\; \theta_{C}\big(d_F^\top y - z_F(r)\big)
    \;\le\; \theta\big(d_F^\top y - z_F(r)\big),
\]
so $d_L^\top y + \theta\,(z_F(r) - d_F^\top y) \le \phi(x,h,\xi)$ for every
$y \in \mathcal{Y}$, and $y = y^{*}$ attains equality.

When the circuits have a combinatorial description, $\theta_{C}$ can be
computed exactly. In the facility location problem of
Section~\ref{subsec:exp-location}, each reservation row contains a single
variable and imposes no condition on circuits, so the circuits are the
alternating paths and even cycles of the bipartite multigraph of stores and
customers, with one arc for reserved and one for on-site purchases. Along a
cycle every store keeps its total, so $d_L^\top g = 0$; along a path the sales
of B change by at most one unit. Hence $\theta_{C} = v/\Delta$, where $\Delta$
is the smallest positive increase of the follower's cost along a path that
lowers the sales of B by one unit. The argument in
Section~\ref{subsec:exp-location} bounds $\Delta$ below by the resolution of the
costs; enumerating the paths instead gives $\Delta$ itself.
\end{remark}

% 실험 절 (INFORMS-IJOC-Template.tex 6.1, "The constants of Section ... are set as follows") 교체 후보 문장:
% The recourse matrix of~\eqref{eq:loc-Q} is totally unimodular, and Remark~\ref{S-rem:theta-circuit} reduces the
% multiplier bound to the smallest increase $\Delta$ of the follower's cost along an alternating path of stores and
% customers that lowers the sales of B. Enumerating the 79{,}140 such paths of the base instance takes 0.1 seconds
% and gives $\Delta = 5$, so $\theta^{U} = v/\Delta = 1$; bounding $\Delta$ below by the resolution of the costs
% would give $\theta^{U} = 5$. The constant matters for computation, because $\theta^{U}$ scales the difference
% of two large terms in $F$: at $x = \{1,2\}$, a 600-second solve of the ISP leaves a relative gap of 95\% with
% $\theta^{U} = 5$ and of 0.2\% with $\theta^{U} = 1$.
% 수치 근거 (x={1,2}, quota 150, δ 0.1, 600 s): Gurobi Ω gap θ=5: 95% (4 스레드) / 15% (전체 스레드, 900 s), θ=1: 0.22% (전체 스레드).
%   α-B&B (worker 12): θ=5 204%, θ=1 28%. logs/diag_abb_base_x12_q150_d0p1_theta*.log
